% ======================================================================
%  KAVRAM: Tool Calling  (Tier 1)  — kompakt
% ======================================================================

\concept{Actions · Eylem}{Tool Calling}{Modelin kendi kafasından çıkıp dış dünyada eylem yapması}

\stage{1}{Problem}
Bir LLM yalnızca metin üretebilir. Bugünün hava durumunu bilemez, veritabanını
sorgulayamaz, e-posta gönderemez, güvenilir aritmetik yapamaz. Ne kadar zeki
olursa olsun, kendi kafasının içinde hapistir --- ve bilmediği bir şeyi
uydurmaya meyillidir.

\keyline{Model ``bilmek''te değil, ``yapmak''ta zayıftır. Eylem için dış dünyaya bir köprü gerekir.}

\stage{2}{Neden Var}
Gerçek görevler eylem ister: veri çekmek, hesaplamak, bir kaydı güncellemek.
Tool calling bu köprüdür: modele kullanabileceği fonksiyonların tanımını
verirsin; model hangisini, hangi argümanlarla çağıracağına karar verip
\emph{yapılandırılmış} bir istek üretir; senin kodun çalıştırıp sonucu modele
geri verir.

\stage{3}{Tarihçe}
\begin{timeline}
\yr{2023} ``plugins'': model dış servis çağırır \tlsep \yr{2023} function calling: yapılandırılmış, şema tabanlı çağrı \tlsep \yr{2024} paralel + çok adımlı tool use standartlaşır \tlsep \yr{2024–26} tool calling, agent'ların temel yapı taşı olur
\end{timeline}

\stage{4}{Tanım}
\begin{problembox}
\textbf{Tool Calling (Function Calling):} Modele araç tanımları (ad, açıklama,
parametre şeması) verilir; model bir aracı \emph{yapılandırılmış} biçimde çağırır
(ör. \texttt{get\_order(id=123)}); senin kodun çalıştırır ve sonucu modele
döndürür. Model çağrıyı \emph{üretir}, çalıştırmaz.
\end{problembox}

\stage{5}{Zihinsel Model}
\begin{mentalbox}
Model, işleri kendi eliyle yapamayan ama \textbf{kusursuz iş emri yazabilen bir
yönetici} gibidir: ``\texttt{get\_weather(city=İstanbul)} çağır.'' Senin kodun
ise emri uygulayıp rapor eden \textbf{ekiptir}. Model \emph{neyi} ve \emph{ne
zaman} yapacağına karar verir; \emph{yapma} işini koda bırakır.
\end{mentalbox}

\stage{6}{Mimari Diyagram}
Tool calling bir döngüdür: model karar verir, kod çalıştırır, sonuç modele döner.
\vspace{3pt}
\begin{center}
\begin{tikzpicture}[font=\sffamily\scriptsize,
   io/.style={draw=softgray, fill=white, rounded corners=8pt, align=center,
              minimum height=9mm, minimum width=18mm, inner sep=3pt},
   m/.style={draw=accentdk, fill=accent!12, rounded corners=2pt, align=center,
             minimum height=11mm, minimum width=20mm},
   k/.style={draw=accent, fill=wash, rounded corners=2pt, align=center,
             minimum height=11mm, minimum width=22mm},
   flow/.style={-{Stealth[length=2mm]}, accent, line width=0.8pt}]
  \node[io] (u) {Kullanıcı};
  \node[m, right=12mm of u] (model) {LLM\\\textcolor{softgray}{karar}};
  \node[k, right=16mm of model] (tool) {Senin kodun\\\textcolor{softgray}{aracı çalıştır}};
  \node[io, below=9mm of model, fill=good!10, draw=good] (ans) {Cevap};

  \draw[flow] (u) -- (model);
  \draw[flow] (model.north) .. controls +(0,6mm) and +(0,6mm) ..
        node[above, font=\sffamily\tiny, text=accentdk]{\texttt{call(args)}} (tool.north);
  \draw[flow] (tool.south) .. controls +(0,-6mm) and +(0,-6mm) ..
        node[below, font=\sffamily\tiny, text=accentdk]{sonuç} (model.south east);
  \draw[flow] (model) -- (ans);
  \node[font=\sffamily\tiny, text=softgray, above=1mm of model] {gerekirse tekrarla};
\end{tikzpicture}
\end{center}
\keyline{Model çağrıyı önerir; güvenlik ve doğrulama \emph{senin kodundadır}. Modelin çıktısı bir niyet, bir komut değildir.}

\stage{7}{Karar Çerçevesi}
Bir yeteneği nasıl vermeli --- prompt'a gömerek mi, RAG ile mi, araç olarak mı?\par
\vspace{2pt}
\noindent\begin{tabularx}{\linewidth}{@{}>{\raggedright\arraybackslash}p{34mm} L@{}}
\arrayrulecolor{hairline}\toprule
{\sffamily\bfseries\small İhtiyaç} & {\sffamily\bfseries\small Doğru araç} \\
\midrule
Canlı / değişken veri (sipariş, hava) & \textbf{Tool calling} \\
Kesin işlem (hesap, kayıt güncelleme) & \textbf{Tool calling} \\
Statik bilgi yığını (dokümanlar) & RAG \\
Sabit kural / küçük bağlam & Prompt \\
\bottomrule
\end{tabularx}
\keyline{İyi araç: tek ve net işlevli, iyi açıklanmış, dar parametreli. Belirsiz/çok geniş araç, modeli yanlış seçime iter.}

\stage{8}{Üretim Senaryosu}
\begin{prodbox}[Üretim Senaryosu · Sipariş Durumu Soran Bot]
\textbf{Bağlam.} Destek botu ``Siparişim nerede?'' sorusuna eskiden uyduruk
tahminler veriyordu --- canlı veriye erişimi yoktu.

\smallskip
\textbf{Çözüm.} \texttt{get\_order\_status(order\_id)} aracı tanımlanır. Model
müşteri kimliğinden sipariş no'yu çıkarır, aracı çağırır, gerçek durumu alır ve
doğal dile çevirir.

\smallskip
\textbf{Sonuç.} Halüsinasyon yerine gerçek, güncel cevap. Yetki kontrolü araçta
yapıldığı için müşteri yalnızca kendi siparişini görebiliyor.
\end{prodbox}

\stage{9}{Anti-pattern}
\begin{antibox}[Anti-pattern · ``50 Araç + Muğlak Açıklama'']
Modele onlarca, üst üste binen ve belirsiz açıklamalı araç vermek; ya da
modelin kendi başına aritmetik/arama yapmasına güvenmek. Model yanlış aracı
seçer, argümanları karıştırır, sessizce hata üretir.

\smallskip
\textbf{Belirti:} ``Model yanlış aracı çağırıyor veya hiç çağırmıyor.''\\
\textbf{Kök neden:} Çok fazla / örtüşen / kötü tanımlı araç.\\
\textbf{Düzeltme:} Araç sayısını azalt, her birini tek işe indir, açıklama ve parametre şemasını netleştir, sonucu doğrula.
\end{antibox}

\stage{10}{Alternatifler \& Ödünleşimler}
\vspace{2pt}
\begin{tabularx}{\linewidth}{@{}>{\raggedright\arraybackslash}p{28mm} L L@{}}
\arrayrulecolor{hairline}\toprule
{\sffamily\bfseries\small Yaklaşım} & {\sffamily\bfseries\small Güçlü} & {\sffamily\bfseries\small Ödünleşim} \\
\midrule
\textbf{Tool calling} & Canlı veri + kesin eylem & Araç tasarımı + güvenlik yükü \\
\addlinespace[2pt]
\textbf{Yalnız prompt} & Basit, sabit görevler & Eylem/canlı veri yapamaz \\
\addlinespace[2pt]
\textbf{RAG} & Statik bilgi getirme & Eylem değil, yalnız okuma \\
\addlinespace[2pt]
\textbf{Sabit kod akışı} & Tam belirli iş & Esnek değil; model kararını kullanmaz \\
\bottomrule
\end{tabularx}

\stage{11}{İlişkili Kavramlar}
\term{MCP} (araçları standartlaştırır) \tlsep \term{Agent} (araç çağrımını döngüye sokar) \tlsep
\term{Function Schema} (parametre sözleşmesi) \tlsep \term{Structured Output} (çağrının biçimi) \tlsep
\term{Authorization} (araç içinde yetki)

\stage{12}{Özet}
\begin{itemize}
  \item Tool calling, modeli ``bilmek''ten ``yapmak''a taşır.
  \item Model çağrıyı \emph{üretir}; çalıştırma, doğrulama ve yetki senin kodundadır.
  \item Canlı veri ve kesin işlem için araç; statik bilgi için RAG; sabit kural için prompt.
  \item İyi araç: tek işlevli, net açıklamalı, dar parametreli.
  \item Az ve net araç, çok ve muğlak araçtan iyidir.
\end{itemize}

\stage{13}{Mühendislik İlkesi}
\begin{principlebox}
\textbf{İlke \#6 — Modelin işi karar vermek; eylemi koda yaptır.}\\[4pt]
\normalsize Modelin çağrısı bir \emph{niyettir}, bir yetki değil. Güvenliği,
doğrulamayı ve sınırı her zaman aracın içinde --- modelin dışında --- uygula.
\end{principlebox}
